\section{Related Work}
\label{sec:related}

\para{LLM sorting and ranking benchmarks.}
\citet{herbold2025sortbench} introduce \textsc{SortBench}, a benchmark that evaluates LLMs on sorting lists of integers, floats, English words, and random strings at lengths from 2 to 256 items.
They find that frontier models achieve high correctness on standard sorting but that reasoning models can \emph{overthink} adversarial tasks---for instance, sorting number words (``three,'' ``seven'') lexicographically rather than numerically.
Our work extends this line of inquiry from syntactic sorting to \emph{semantic} sorting: ordering items by factual properties like weight, elevation, or chronological date.

\para{Sorting algorithms for LLMs.}
\citet{zhao2025access} systematically compare physical implementations for LLM-based ordering---pointwise scoring, pairwise comparison, and listwise ranking---on factual tasks (NBA player heights, country populations) and semantic tasks (passage relevance).
They find that pointwise methods achieve high accuracy on factual tasks because LLMs can directly recall memorized values, while comparison-based methods excel on semantic tasks.
Our work complements theirs by holding the algorithm fixed (one-pass listwise) and varying the \emph{property}, revealing which types of factual knowledge are most accessible.

\para{Listwise ranking with LLMs.}
\citet{sun2023chatgpt} propose \textsc{RankGPT}, which uses LLMs as zero-shot passage re-rankers via a sliding-window listwise approach.
\citet{tang2023found} address positional bias in listwise ranking by shuffling input order across multiple runs and aggregating results, achieving substantial gains for smaller models.
\citet{zeng2024rankfusion} identify order inconsistency and transitive inconsistency as fundamental challenges in LLM-based ranking and propose fusion strategies to mitigate them.
These works focus on passage relevance ranking; we focus on factual ordering where ground truth is objective and measurable.

\para{Probing factual knowledge in LLMs.}
A growing body of work probes what LLMs ``know'' by testing factual recall~\cite{suzgun2022challenging}.
\citet{patel2025semantic} formalize semantic operators---including a sort operator---for LLM-based data processing, demonstrating that LLMs can perform rich data transformations.
Our ordering-based probe complements cloze-style fact probes by testing \emph{relational} knowledge (A $>$ B) rather than pointwise recall.

Unlike prior work, we systematically vary the \emph{type} of property being ordered and measure accuracy across a taxonomy of property categories.
This enables us to identify a hierarchy of feature accessibility that generalizes across models and provides practical guidance for when to trust LLM-based ordering.
